% Generated draft: requires mathematical and source review.
\documentclass[11pt]{article}
\usepackage[margin=25mm]{geometry}
\usepackage{fontspec}
\setmainfont{Latin Modern Roman}
\usepackage{amsmath,amssymb,mathtools,unicode-math}
\setmathfont{Latin Modern Math}
\usepackage{enumitem,needspace,xurl,hyperref}
\setlength{\parindent}{0pt}
\setlength{\parskip}{5pt}
\setlength{\emergencystretch}{3em}
\newcommand{\N}{\mathbb{N}}
\newcommand{\Z}{\mathbb{Z}}
\newcommand{\Q}{\mathbb{Q}}
\newcommand{\R}{\mathbb{R}}
\newcommand{\C}{\mathbb{C}}
\newcommand{\F}{\mathbb{F}}
\newcommand{\A}{\mathbb{A}}
\newcommand{\PP}{\mathbb{P}}
\newcommand{\E}{\mathbb{E}}
\newcommand{\eps}{\varepsilon}
\newcommand{\dd}{\,\mathrm{d}}
\begin{document}
\section*{UnsolvedMath ID: 30003687}
\textbf{Source problem code:} catalogue entry
\subsection{Definitions and mathematical statement}
For reference, a standard classical notion is uniform distribution modulo \(1\). Let \((x_n)_{n\geq 1}\) be a sequence in \(\R\), and write \(\{x\}=x-\lfloor x\rfloor\). The sequence is uniformly distributed modulo \(1\) if
\[
\forall\,0\leq a<b\leq 1,\qquad
\lim_{N\to\infty}\frac{1}{N}\#\bigl\{1\leq n\leq N:a\leq\{x_n\}<b\bigr\}=b-a.
\]
This definition supplies a conventional reference point; the source does not specify whether distribution modulo \(1\) is the intended setting.

\textbf{Source question.} How is algorithmic randomness related to the classical theory of uniform distribution of sequences of real numbers?

\textbf{Unresolved formulation.} The supplied question is a research direction rather than a quantified conjecture. It specifies neither a particular notion of algorithmic randomness nor an underlying probability space or representation of random objects. It also does not specify how such objects determine real sequences, which sequences or transformations are permitted, or which implication, equivalence, or separation is sought. Consequently, the record does not determine a unique precise mathematical assertion. Choosing these ingredients would require additional source evidence.
\subsection{Short English statement}
The question asks for relationships between algorithmic randomness and classical uniform distribution. A standard definition of uniform distribution modulo \(1\) is provided as context, but the intended randomness notion, sequence construction, distribution setting, and exact relationship remain unspecified. The supplied record labels the problem partially solved; this is a supplied assessment requiring review.
\subsection{Sources}
\begin{itemize}
\item Supplied local UnsolvedMath registry record, permanent ID \(30003687\), source problem code \texttt{catalogue entry}, title \emph{Algorithmic Randomness and Uniform Distribution}, in the Oberwolfach Reports Open Problems collection. The project identifies \texttt{data/unsolvedmath/problems.json} and \texttt{catalog/problems.json} as local registries; their contents have not been independently inspected for this draft.
\item Supplied report citation: \emph{Computability Theory (2019)}, Oberwolfach Reports, DOI \texttt{10.4171/owr/2018/1}, pp.\ 5--41. \url{https://doi.org/10.4171/owr/2018/1}. The background describes Oberwolfach Reports as documenting mathematical workshops and their open-problem discussions. The record gives proposed year \(2018\) and supplies no proposer attribution; the relationship between that year and the citation's \(2019\) designation is unresolved.
\item Supplied literature reference: Avigad, \emph{Uniform distribution and algorithmic randomness}. \url{https://arxiv.org/abs/1203.4126}.
\item Supplied literature assessment, checked \texttt{2026-08-21}: status \emph{partially solved}. The supplied curation attributes several implications and separations between algorithmic randomness notions and uniform distribution to Avigad, while stating that the broad relationship asked about is not exhausted. Its listed sources are the Avigad reference and Oberwolfach report above. This is supplied dated curation derived from the OWR record, not an independently verified literature assessment; time-sensitive status and cited claims require independent verification.
\end{itemize}
\end{document}
